\begingroup
\renewcommand{\clearpageforsection}{}
\exercisesheader{}
\endgroup


% 17

\eoce{\qt{Eksperimen sosial, Bagian I\label{social_experiment_conditions}} Sebuah ``eksperimen 
sosial" yang dilakukan oleh sebuah program TV menyelidiki tindakan orang ketika melihat 
seorang perempuan dengan memar yang sangat kentara dirundung oleh pacarnya. Pada dua 
kesempatan berbeda di restoran yang sama, pasangan yang sama diperagakan. Dalam 
satu skenario, perempuan tersebut berpakaian ``provokatif'', sedangkan dalam 
skenario lainnya ia berpakaian ``konservatif''. Tabel di bawah menunjukkan berapa 
banyak pengunjung restoran yang hadir dalam setiap skenario dan apakah 
mereka turun tangan.
\begin{center}
\begin{tabular}{ll cc c} 
            &               & \multicolumn{2}{c}{\textit{Skenario}} \\
\cline{3-4}
                            &       & Provokatif   & Konservatif  & Jumlah \\
\cline{2-5}
\multirow{2}{*}{\textit{Turun tangan}} &Ya & 5   & 15  & 20    \\
                            &Tidak     & 15      & 10  & 25 \\
\cline{2-5}
                            &Jumlah  & 20      & 25  & 45 \\
\end{tabular}
\end{center}
Jelaskan mengapa distribusi sampling selisih antara 
proporsi orang yang turun tangan dalam skenario provokatif dan konservatif 
tidak mengikuti distribusi yang mendekati normal.
}{}

% 18

\eoce{\qt{Keberhasilan transplantasi jantung\label{heart_transplant_conditions}} Studi 
Transplantasi Jantung Universitas Stanford dilakukan untuk menentukan apakah suatu 
program transplantasi jantung eksperimental memperpanjang masa hidup. Setiap pasien 
yang memasuki program tersebut secara resmi ditetapkan sebagai calon transplantasi jantung, 
yang berarti pasien itu sakit parah dan mungkin memperoleh manfaat dari jantung baru. Pasien 
ditempatkan secara acak ke dalam kelompok perlakuan dan kelompok kontrol. Pasien dalam 
kelompok perlakuan menerima transplantasi, sedangkan pasien dalam kelompok kontrol 
tidak menerimanya. Tabel di bawah menunjukkan jumlah pasien yang bertahan hidup dan meninggal dalam setiap 
kelompok. \footfullcite{Turnbull+Brown+Hu:1974}\vspace{-2mm}
\begin{center}
\begin{tabular}{rcc}
\hline
            & kontrol   & perlakuan \\ 
\hline
hidup       & 4         & 24 \\ 
meninggal   & 30        & 45 \\ 
\hline
\end{tabular}
\end{center}
Misalkan kita ingin mengestimasi selisih tingkat kelangsungan hidup antara 
kelompok kontrol dan kelompok perlakuan dengan menggunakan interval kepercayaan.
Jelaskan mengapa kita tidak dapat menyusun interval semacam itu dengan menggunakan 
aproksimasi normal. Apa yang mungkin keliru jika kita tetap menyusun interval kepercayaan 
meskipun terdapat masalah ini?
}{}

% 19

\eoce{\qt{Jenis kelamin dan preferensi warna\label{gender_color_preference_CI_concept}} 
Sebuah studi menanyakan kepada 1,924 mahasiswa sarjana laki-laki dan 3,666 mahasiswa sarjana perempuan
warna kesukaan mereka.
Interval kepercayaan 95\% untuk selisih antara 
proporsi laki-laki dan perempuan yang warna kesukaannya hitam 
$(p_{male} - p_{female})$ dihitung sebesar (0.02, 0.06).
Berdasarkan 
informasi ini, tentukan apakah pernyataan-pernyataan berikut tentang
mahasiswa sarjana benar atau salah, dan 
jelaskan alasan Anda untuk setiap pernyataan yang Anda nyatakan salah.
\footfullcite{Ellis:2001}
\begin{parts}
\item Kita yakin 95\% bahwa proporsi sebenarnya laki-laki yang warna 
kesukaannya hitam berada antara 2\% lebih rendah dan 6\% lebih tinggi daripada proporsi sebenarnya perempuan 
yang warna kesukaannya hitam.
\item Kita yakin 95\% bahwa proporsi sebenarnya laki-laki yang warna 
kesukaannya hitam 2\% hingga 6\% lebih tinggi daripada proporsi sebenarnya perempuan yang 
warna kesukaannya hitam.
\item Sebanyak 95\% sampel acak akan menghasilkan interval kepercayaan 95\% yang 
memuat selisih sebenarnya antara proporsi populasi laki-laki dan 
perempuan yang warna kesukaannya hitam.
\item Kita dapat menyimpulkan bahwa terdapat selisih yang signifikan antara 
proporsi laki-laki dan perempuan yang warna kesukaannya hitam, serta bahwa 
selisih antara kedua proporsi sampel terlalu besar untuk secara masuk akal 
disebabkan oleh kebetulan.
\item Interval kepercayaan 95\% untuk $(p_{female} - p_{male})$ tidak dapat 
dihitung hanya dengan informasi yang diberikan dalam latihan ini.
\end{parts}
}{}

% Boundary-clean reader removes a legacy forced page break.

% 20

\eoce{\qt{Penutupan pemerintah\label{government_shutdown_CI_concept}}
Penutupan pemerintah federal Amerika Serikat tahun 2018--2019 berlangsung
dari 22 Desember 2018 hingga 25 Januari 2019, selama 35 hari.
Jajak pendapat Survey USA terhadap sampel acak 614 warga Amerika pada
periode tersebut melaporkan bahwa 48\% orang yang berpenghasilan kurang dari \$40,000 per
tahun dan 55\% orang yang berpenghasilan \$40,000 atau lebih per tahun mengatakan bahwa
penutupan pemerintah sama sekali tidak memengaruhi mereka secara pribadi.
Interval kepercayaan 95\% untuk $(p_\text{$<$40K} - p_\text{$\ge$40K})$,
dengan $p$ sebagai proporsi orang yang mengatakan bahwa penutupan pemerintah
sama sekali tidak memengaruhi mereka secara pribadi, adalah (-0.16, 0.02).
Berdasarkan informasi ini, tentukan apakah pernyataan-pernyataan berikut 
benar atau salah, dan jelaskan alasan Anda jika menyatakan suatu pernyataan
salah.\footfullcite{data:govt_shuthown}
\begin{parts}
\item
    Pada tingkat signifikansi 5\%, data memberikan
    bukti yang meyakinkan tentang adanya selisih nyata dalam proporsi orang yang
    tidak terdampak secara pribadi antara warga Amerika yang berpenghasilan kurang dari
    \$40,000 per tahun dan warga Amerika yang berpenghasilan \$40,000 atau lebih per tahun.
\item
    Kita yakin 95\% bahwa terdapat antara 16\% lebih banyak dan 2\% lebih sedikit warga Amerika
    yang berpenghasilan kurang dari \$40,000 per tahun dan sama sekali tidak terdampak secara pribadi
    oleh penutupan pemerintah dibandingkan dengan mereka yang berpenghasilan
    \$40,000 atau lebih per tahun.
\item
    Interval kepercayaan 90\% untuk
    $(p_\text{$<$40K} - p_\text{$\ge$40K})$
    akan lebih lebar daripada interval $(-0.16, 0.02)$.
\item
    Interval kepercayaan 95\% untuk
    $(p_\text{$\ge$40K} - p_\text{$<$40K})$
    adalah (-0.02, 0.16).
\end{parts}
% p1 = 0.48
% p2 = 0.55
% n1 = 162
% n2 = 452
% ((p1 - p2) + c(-1,1) * 1.96 * sqrt( (p1*(1-p1)/n1) + (p2*(1-p2)/n2)) ) %>% round(2)
% (-0.16, 0.02)
}{}

% 21

\eoce{\qt{Rencana Kesehatan Nasional,
    Bagian III\label{national_health_plan_CI_replaced}}
Latihan~\ref{national_health_plan_HT}
menyajikan hasil jajak pendapat yang mengevaluasi dukungan terhadap
``Rencana Kesehatan Nasional'' berlabel umum
di Amerika Serikat.
Sebanyak 79\% dari 347 Demokrat dan 55\% dari 617 Independen
mendukung Rencana Kesehatan Nasional.
\begin{parts}
\item
    Hitung interval kepercayaan 95\% untuk
    selisih antara proporsi Demokrat
    dan Independen yang mendukung Rencana
    Kesehatan Nasional $(p_{D} - p_{I})$, lalu interpretasikan
    interval tersebut dalam konteks ini.
    Kami telah memeriksa syarat-syaratnya untuk Anda.
\item
    Benar atau salah:
    Jika pada saat jajak pendapat ini kita memilih seorang Demokrat acak dan seorang
    Independen acak, Demokrat tersebut lebih
    mungkin mendukung Rencana Kesehatan
    Nasional daripada Independen tersebut.
\end{parts}
}{}

% 22

\eoce{\qt{Kurang tidur, CA vs. OR, Bagian I\label{sleep_OR_CA_CI}} Menurut 
laporan tentang kurang tidur oleh Pusat Pengendalian dan Pencegahan Penyakit, 
proporsi penduduk California yang melaporkan tidak cukup beristirahat atau tidur 
selama masing-masing dari 30 hari sebelumnya adalah 8.0\%, sedangkan proporsi ini sebesar 8.8\% 
untuk penduduk Oregon. Data tersebut didasarkan pada sampel acak sederhana yang terdiri atas 11,545 
penduduk California dan 4,691 penduduk Oregon. Hitung interval kepercayaan 95\% 
untuk selisih antara proporsi warga California dan Oregon yang 
kurang tidur, lalu interpretasikan interval tersebut dalam konteks data.\footfullcite{data:sleepCAandOR}
}{}

% 23

\eoce{\qt{Pengeboran lepas pantai, Bagian I\label{offshore_drill_edu_dontknow_HT}}
Sebuah survei bertanya kepada 827 pemilih terdaftar di California yang diambil secara acak,
``Apakah Anda mendukung? Atau apakah Anda menentang? Pengeboran minyak dan gas alam
di lepas Pantai California? Atau apakah Anda belum cukup tahu untuk menjawab?''
Di bawah ini ditampilkan distribusi 
respons, yang dipisahkan berdasarkan apakah responden lulus 
perguruan tinggi atau tidak. \footfullcite{data:prop19_and_offshoreDrill} \\[1.3mm]

\noindent\begin{minipage}[c]{0.6\textwidth}
\begin{parts}
\item Berapa persen lulusan perguruan tinggi dan berapa persen bukan lulusan 
perguruan tinggi dalam sampel ini yang belum cukup tahu untuk berpendapat mengenai pengeboran 
minyak dan gas alam di lepas Pantai California?
\item Lakukan uji hipotesis untuk menentukan apakah data tersebut memberikan bukti kuat 
bahwa proporsi lulusan perguruan tinggi yang tidak berpendapat 
mengenai persoalan ini berbeda dari proporsi bukan lulusan perguruan tinggi.
\end{parts}
\end{minipage}
\begin{minipage}[c]{0.4\textwidth}
\begin{center}
\begin{tabular}{l c c}
			& \multicolumn{2}{c}{\textit{Lulusan perguruan tinggi}} \\
\cline{2-3}
			& Ya		& Tidak	\\
\cline{1-3}
Mendukung		& 154		& 132	\\
Menentang		& 180		& 126	\\
Tidak tahu	& 104		& 131	\\
\cline{1-3}
  Jumlah		& 438		& 389		
\end{tabular}
\end{center}
\end{minipage}
}{}

% 24
 
\eoce{\qt{Kurang tidur, CA vs. OR, Bagian II\label{sleep_OR_CA_HT}} 
Latihan~\ref{sleep_OR_CA_CI} menyajikan data tentang tingkat kurang tidur di 
California dan Oregon. Proporsi penduduk California yang 
melaporkan tidak memperoleh istirahat atau tidur yang cukup pada setiap hari selama 30 hari sebelumnya adalah 
8.0\%, sedangkan proporsi tersebut adalah 8.8\% untuk penduduk Oregon. Data ini 
didasarkan pada sampel acak sederhana yang terdiri atas 11,545 penduduk California dan 4,691 penduduk Oregon. 

\begin{parts}
\item Lakukan uji hipotesis untuk menentukan apakah data ini memberikan bukti kuat 
bahwa tingkat kurang tidur berbeda di kedua negara bagian tersebut. 
(Pengingat: Periksa syarat-syaratnya.)
\item Kesimpulan uji pada bagian (a) mungkin keliru. Jika 
demikian, jenis galat apa yang terjadi?
\end{parts}
}{}

% Boundary-clean reader removes a legacy forced page break.

% 25

\eoce{\qt{Pengeboran lepas pantai, Bagian II\label{offshore_drill_edu_support_HT}}
Hasil jajak pendapat yang mengevaluasi dukungan terhadap pengeboran minyak
dan gas alam di lepas pantai California diperkenalkan
dalam Latihan~\ref{offshore_drill_edu_dontknow_HT}.
\begin{center}
\begin{tabular}{l c c}
				& \multicolumn{2}{c}{\textit{Lulusan Perguruan Tinggi}} \\
\cline{2-3}
						& Ya		& Tidak				\\
\cline{1-3}
Mendukung		& 154		& 132			\\
Menentang		& 180		& 126			\\
Tidak tahu	& 104		& 131			\\
\cline{1-3}
 Total		& 438		& 389		
\end{tabular}
\end{center}
\begin{parts}
\item Berapa persen lulusan perguruan tinggi dan berapa persen responden yang bukan 
lulusan perguruan tinggi dalam sampel ini yang mendukung pengeboran minyak dan gas alam di lepas pantai 
California?
\item Lakukan uji hipotesis untuk menentukan apakah data memberikan bukti kuat 
bahwa proporsi lulusan perguruan tinggi yang mendukung pengeboran lepas pantai di California 
berbeda dari proporsi responden yang bukan lulusan perguruan tinggi.
\end{parts}
}{}

% 26

\eoce{\qt{Pemindaian seluruh tubuh, Bagian I\label{full_body_scan_HT_Error}} Sebuah artikel berita 
melaporkan bahwa ``warga Amerika memiliki pandangan yang berbeda tentang dua praktik yang mungkin merepotkan 
dan invasif yang dapat diterapkan bandara untuk mengungkap kemungkinan 
serangan teroris.'' Berita ini didasarkan pada survei terhadap 
sampel acak 1,137 orang dewasa di seluruh negeri. Salah satu pertanyaan dalam 
survei tersebut berbunyi, ``Sejumlah bandara kini 
menggunakan mesin sinar-X digital `seluruh tubuh' untuk memindai penumpang secara elektronik 
di antrean pemeriksaan keamanan bandara. Menurut Anda, apakah mesin sinar-X baru ini 
seharusnya digunakan atau tidak digunakan di bandara?'' Berikut ringkasan tanggapan 
berdasarkan afiliasi partai. \footfullcite{news:fullBodyScan}
\begin{center}
\begin{tabular}{ll  cc c} 
            &   & \multicolumn{3}{c}{\textit{Afiliasi Partai}} \\
\cline{3-5}
                                &           & Republikan & Demokrat & Independen   \\
\cline{2-5}
\multirow{3}{*}{\textit{Jawaban}}& Seharusnya    & 264        & 299      & 351 \\
                                & Seharusnya tidak& 38         & 55       & 77 \\
                                & Tidak tahu/Tanpa jawaban & 16 & 15    & 22 \\
\cline{2-5}
                                & Total      & 318       & 369      & 450
\end{tabular}
\end{center}
\begin{parts}
\item Lakukan uji hipotesis yang sesuai untuk mengevaluasi apakah terdapat 
perbedaan antara proporsi Republikan dan Demokrat yang menilai bahwa pemindaian seluruh 
tubuh seharusnya diterapkan di bandara. Anggap semua syarat yang relevan 
terpenuhi.
\item Kesimpulan uji pada bagian (a) mungkin keliru, yang berarti terjadi 
galat pengujian. Jika terjadi galat, apakah itu galat Tipe~1 atau galat Tipe~2? 
Jelaskan.
\end{parts}
}{}

% 27

\eoce{\qt{Pekerja transportasi yang kurang tidur\label{sleep_deprived_driver_HT}} 
National Sleep Foundation melakukan survei tentang kebiasaan tidur 
pekerja transportasi yang dipilih secara acak dan sampel kontrol pekerja di luar sektor transportasi. 
Hasil survei tersebut ditampilkan di bawah ini.
\footfullcite{data:sleepTransport}\vspace{-1.8mm}
\begin{center}
\begin{tabular}{l c c c c c }
						& 			& \multicolumn{4}{c}{\textit{Profesional Transportasi}} \\
\cline{3-6}
			& 			& 		& Pengemudi	& Operator		& Pengemudi Bus/ 		\\
			& \textit{Kontrol}& Pilot	& Truk	& Kereta	& Taksi/Limusin	\\
\cline{1-6}
Tidur kurang dari 6 jam	& 35		& 19		& 35	& 29	& 21	\\
Tidur 6 hingga 8 jam		& 193		& 132	    & 117	& 119	& 131	\\
Tidur lebih dari 8 jam			& 64		& 51		& 51	& 32	& 58	\\
\cline{1-6}
Total						& 292		& 202	    & 203	& 180	& 210		
\end{tabular}
\end{center}\vspace{-1.2mm}
Lakukan uji hipotesis untuk mengevaluasi apakah data ini memberikan bukti adanya 
perbedaan antara proporsi pengemudi truk dan pekerja di luar sektor transportasi 
(kelompok kontrol) yang tidur kurang dari 6 jam per hari, yakni 
yang dianggap kurang tidur.
}{}

% Boundary-clean reader removes a legacy forced page break.

% 28

\eoce{\qt{Vitamin prenatal dan autisme\label{prenatal_vitamin_autism_HT}} 
Peneliti yang mempelajari kaitan antara penggunaan vitamin prenatal dan autisme 
menyurvei para ibu dari suatu sampel acak anak berusia 24--60 bulan dengan 
autisme dan mengambil sampel acak terpisah dari anak dengan perkembangan 
tipikal. Tabel di bawah menunjukkan jumlah ibu pada setiap kelompok yang 
menggunakan dan tidak menggunakan vitamin prenatal selama tiga bulan sebelum 
kehamilan (periode perikonsepsi).\footfullcite{Schmidt:2011}\vspace{-1.8mm}
\begin{center}
\begin{tabular}{l l c c c}
		&			& \multicolumn{2}{c}{\textit{Autisme}}	&		\\
\cline{3-4}
		&			& Autisme		& Perkembangan tipikal		& Total	\\
\cline{2-5}
\textit{Vitamin prenatal}	& Tanpa vitamin	& 111	& 70		& 181	\\
\textit{perikonsepsi}	& Dengan vitamin	& 143		& 159		& 302	\\
\cline{2-5}
							& Total		& 254		& 229		& 483
\end{tabular}
\end{center}\vspace{-4.2mm}
\begin{parts}
\item Nyatakan hipotesis yang sesuai untuk menguji independensi antara penggunaan 
vitamin prenatal selama tiga bulan sebelum kehamilan dan autisme.
\item Selesaikan uji hipotesis dan nyatakan kesimpulan yang sesuai. 
(Pengingat: Verifikasi semua syarat yang diperlukan untuk uji tersebut.)
\item Sebuah artikel New York Times yang melaporkan penelitian ini berjudul ``Vitamin 
Prenatal Mungkin Mencegah Autisme''. Menurut Anda, apakah judul artikel tersebut 
tepat? Jelaskan jawaban Anda. Selain itu, usulkan judul alternatif.
\footfullcite{news:prenatalVitAutism}
\end{parts}
}{}

% 29

\eoce{\qt{HIV di Afrika sub-Sahara\label{hiv_africa_HT}}
Pada Juli 2008, US National Institutes of Health mengumumkan
bahwa sebuah studi klinis dihentikan lebih awal karena hasil yang tidak terduga.
Populasi studi terdiri atas perempuan terinfeksi HIV di Afrika
sub-Sahara yang telah diberi satu dosis nevirapine (obat untuk HIV)
saat melahirkan, untuk mencegah penularan HIV kepada bayi.
Studi tersebut merupakan perbandingan acak atas kelanjutan
pengobatan perempuan (setelah persalinan berhasil)
dengan nevirapine versus lopinavir, obat kedua yang digunakan untuk mengobati HIV.
Sebanyak 240 perempuan berpartisipasi dalam studi;
120 ditempatkan secara acak ke masing-masing dari kedua pengobatan.
Dua puluh empat minggu setelah memulai pengobatan studi,
setiap perempuan diperiksa untuk menentukan apakah infeksi HIV
semakin memburuk (hasil yang disebut \textit{kegagalan virologis}).
Dua puluh enam dari 120 perempuan yang diobati dengan nevirapine mengalami
kegagalan virologis, sedangkan 10 dari 120 perempuan yang diobati dengan
obat lainnya mengalami kegagalan virologis.\footfullcite{Lockman:2007}
\begin{parts}
\item Buat tabel dua arah yang menyajikan hasil studi ini.
\item Nyatakan hipotesis yang sesuai untuk menguji perbedaan tingkat kegagalan virologis
antara kedua kelompok pengobatan.
\item Selesaikan uji hipotesis dan nyatakan kesimpulan yang sesuai. 
(Pengingat: Verifikasi semua syarat yang diperlukan untuk uji tersebut.)
\end{parts}
}{}

% 30

\eoce{\qt{Satu apel sehari menjauhkan dokter
    \label{apple_doctor_HT_concept}}
Seorang guru pendidikan jasmani di sebuah sekolah menengah yang ingin
meningkatkan kesadaran tentang gizi dan kesehatan
bertanya kepada para siswanya pada awal semester
apakah mereka memercayai ungkapan
``satu apel sehari menjauhkan dokter'',
dan 40\% siswa menjawab ya.
Sepanjang semester, ia memulai setiap kelas dengan
diskusi singkat mengenai studi yang menyoroti dampak positif
dari mengonsumsi lebih banyak buah dan sayuran.
Ia melakukan survei tentang satu-apel-sehari yang sama pada akhir
semester, dan kali ini 60\% siswa
menjawab ya.
Dapatkah ia menggunakan metode dua proporsi dari bagian ini
untuk analisis tersebut?
Jelaskan alasan Anda.
}{}
